% GENERATED FILE. Edit canonical lessons, then run npm run generate:books.
% canonical-source-hash: fnv1a64:da7a66ddf86799ea
% canonical-lessons: PA-C240-sil, PA-C240-pihna, PA-C240-ukkhli, PA-C240-mohla, PA-C240-kuttna

\chapter{Grinding and Pounding}
\label{ch:pa-a2-grinding-and-pounding}
\hlchaptermodality{240}

\begin{chapteropening}
\textbf{By the end of this chapter:} \emph{I can name the tools that grind spices.}

Complete the last lesson of chapter 240: I can name the tools that grind spices.
\end{chapteropening}

\section[\texorpdfstring{\pa{ਸਿਲ} (sil)}{ਸਿਲ (sil)}]{\pa{ਸਿਲ} (sil) — a grinding stone}

\label{lesson:PA-C240-sil}



\begin{quote}
Before the new one: say the Punjabi for to cover, then the Punjabi for ghee.
\end{quote}

\subsection*{You'll want to know: \pa{ਸਿਲ}}
\textbf{\pa{ਸਿਲ}} — \emph{sil} — \textquotedblleft{}a grinding stone\textquotedblright{}.

\begin{quote}
\pa{ਸਿਲ} \pa{ਭਾਰੀ} \pa{ਪੱਥਰ} \pa{ਹੈ।}
\end{quote}

\begin{quote}
\emph{sil bhārī patthar hai.}
\end{quote}

\begin{quote}
The grinding stone is a heavy stone.
\end{quote}

A feminine noun.

\begin{cousinweb}
\textbf{\pa{ਸਿਲ}} (\emph{sil}) is inherited from Sanskrit \emph{śilā}, \textquotedblleft{}rock, crag\textquotedblright{}; CDIAL lists Punjabi \emph{sil}, \textquotedblleft{}stone slab used for grinding spices\textquotedblright{} (CDIAL 12459).
\end{cousinweb}

\begin{grammarlens}[title={What you've built}]
One more word for this chapter.
\end{grammarlens}

\subsection*{Your turn}
\begin{itemize}
\raggedright
  \item \emph{Say it:} \emph{sil}
  \item \emph{Say it:} \emph{sil}, once more
  \item \emph{Say it:} \emph{sil bhārī patthar hai.} — The grinding stone is a heavy stone.
  \item \emph{Recall:} say the Punjabi for to cover, then the Punjabi for ghee, then say \emph{sil} again
\end{itemize}

\begin{tcolorbox}[breakable,colback=teal!4,colframe=teal!35!black,title={Before you move on}]
What does \pa{ਸਿਲ} mean? (\textquotedblleft{}A grinding stone\textquotedblright{}.) Say it once more.
\end{tcolorbox}
\section[\texorpdfstring{\pa{ਪੀਹਣਾ} (pīhṇā)}{ਪੀਹਣਾ (pīhṇā)}]{\pa{ਪੀਹਣਾ} (pīhṇā) — to grind}

\label{lesson:PA-C240-pihna}



\begin{quote}
Before the new one: say the Punjabi for ghee, then the Punjabi for a grinding stone.
\end{quote}

\subsection*{You'll want to know: \pa{ਪੀਹਣਾ}}
\textbf{\pa{ਪੀਹਣਾ}} — \emph{pīhṇā} — \textquotedblleft{}to grind\textquotedblright{}.

\begin{quote}
\pa{ਕੱਲ੍ਹ} \pa{ਮੈਂ} \pa{ਸਿਲ} \pa{ਉੱਤੇ} \pa{ਮਸਾਲਾ} \pa{ਪੀਹਾਂਗੀ।}
\end{quote}

\begin{quote}
\emph{kall maiṁ sil utte masālā pīhāṅgī.}
\end{quote}

\begin{quote}
Tomorrow I will grind spices on the stone.
\end{quote}

Past (I …): \pa{ਮੈਂ} \pa{ਮਸਾਲਾ} \pa{ਪੀਹਿਆ}, \emph{maiṁ masālā pīhiā}.

\begin{cousinweb}
\textbf{\pa{ਪੀਹਣਾ}} (\emph{pīhṇā}) is inherited from Sanskrit \emph{piṁṣati}, \textquotedblleft{}grinds\textquotedblright{}; CDIAL lists Punjabi \emph{pīhṇā}, \textquotedblleft{}to grind\textquotedblright{} (CDIAL 8142).
\end{cousinweb}

\begin{grammarlens}[title={What you've built}]
One more word for this chapter.
\end{grammarlens}

\subsection*{Your turn}
\begin{itemize}
\raggedright
  \item \emph{Say it:} \emph{pīhṇā}
  \item \emph{Say it:} \emph{pīhṇā}, once more
  \item \emph{Say it:} \emph{kall maiṁ sil utte masālā pīhāṅgī.} — Tomorrow I will grind spices on the stone.
  \item \emph{Recall:} say the Punjabi for ghee, then the Punjabi for a grinding stone, then say \emph{pīhṇā} again
  \item \emph{Write it:} copy \textbf{\pa{ਪੀਹਣਾ}} once, checking the whole word
\end{itemize}

\begin{tcolorbox}[breakable,colback=teal!4,colframe=teal!35!black,title={Before you move on}]
What does \pa{ਪੀਹਣਾ} mean? (\textquotedblleft{}To grind\textquotedblright{}.) Say it once more.
\end{tcolorbox}
\section[\texorpdfstring{\pa{ਉੱਖਲੀ} (ukkhlī)}{ਉੱਖਲੀ (ukkhlī)}]{\pa{ਉੱਖਲੀ} (ukkhlī) — a wooden mortar}

\label{lesson:PA-C240-ukkhli}



\begin{quote}
Before the new one: say the Punjabi for a grinding stone, then the Punjabi for to grind.
\end{quote}

\subsection*{You'll want to know: \pa{ਉੱਖਲੀ}}
\textbf{\pa{ਉੱਖਲੀ}} — \emph{ukkhlī} — \textquotedblleft{}a wooden mortar\textquotedblright{}.

\begin{quote}
\pa{ਉੱਖਲੀ} \pa{ਵਿੱਚ} \pa{ਹਲਦੀ} \pa{ਹੈ।}
\end{quote}

\begin{quote}
\emph{ukkhlī vich haldī hai.}
\end{quote}

\begin{quote}
There is turmeric in the mortar.
\end{quote}

A feminine noun.

\begin{cousinweb}
\textbf{\pa{ਉੱਖਲੀ}} (\emph{ukkhlī}) is inherited from Sanskrit \emph{ulūkhala}, \textquotedblleft{}a wooden mortar\textquotedblright{}; CDIAL lists Punjabi \emph{ukkhalī}, \textquotedblleft{}small wooden mortar\textquotedblright{} (CDIAL 2360). CDIAL notes that the Sanskrit word itself may come from Dravidian.
\end{cousinweb}

\begin{grammarlens}[title={What you've built}]
One more word for this chapter.
\end{grammarlens}

\subsection*{Your turn}
\begin{itemize}
\raggedright
  \item \emph{Say it:} \emph{ukkhlī}
  \item \emph{Say it:} \emph{ukkhlī}, once more
  \item \emph{Say it:} \emph{ukkhlī vich haldī hai.} — There is turmeric in the mortar.
  \item \emph{Recall:} say the Punjabi for a grinding stone, then the Punjabi for to grind, then say \emph{ukkhlī} again
  \item \emph{Write it:} copy \textbf{\pa{ਉੱਖਲੀ}} once, checking the whole word
\end{itemize}

\begin{tcolorbox}[breakable,colback=teal!4,colframe=teal!35!black,title={Before you move on}]
What does \pa{ਉੱਖਲੀ} mean? (\textquotedblleft{}A wooden mortar\textquotedblright{}.) Say it once more.
\end{tcolorbox}
\section[\texorpdfstring{\pa{ਮੋਹਲਾ} (mohlā)}{ਮੋਹਲਾ (mohlā)}]{\pa{ਮੋਹਲਾ} (mohlā) — a pestle}

\label{lesson:PA-C240-mohla}



\begin{quote}
Before the new one: say the Punjabi for to grind, then the Punjabi for a wooden mortar.
\end{quote}

\subsection*{You'll want to know: \pa{ਮੋਹਲਾ}}
\textbf{\pa{ਮੋਹਲਾ}} — \emph{mohlā} — \textquotedblleft{}a pestle\textquotedblright{}.

\begin{quote}
\pa{ਉੱਖਲੀ} \pa{ਅਤੇ} \pa{ਮੋਹਲਾ} \pa{ਰਸੋਈ} \pa{ਵਿੱਚ} \pa{ਹਨ।}
\end{quote}

\begin{quote}
\emph{ukkhlī ate mohlā rasoī vich han.}
\end{quote}

\begin{quote}
The mortar and pestle are in the kitchen.
\end{quote}

A masculine noun.

\begin{cousinweb}
\textbf{\pa{ਮੋਹਲਾ}} (\emph{mohlā}) is inherited from Sanskrit \emph{musala}, \textquotedblleft{}pestle\textquotedblright{}; CDIAL lists Punjabi \emph{mohlā}, \textquotedblleft{}pestle\textquotedblright{} (CDIAL 10223). CDIAL notes that the Sanskrit word is probably not Indo-Aryan in origin.
\end{cousinweb}

\begin{grammarlens}[title={What you've built}]
One more word for this chapter.
\end{grammarlens}

\subsection*{Your turn}
\begin{itemize}
\raggedright
  \item \emph{Say it:} \emph{mohlā}
  \item \emph{Say it:} \emph{mohlā}, once more
  \item \emph{Say it:} \emph{ukkhlī ate mohlā rasoī vich han.} — The mortar and pestle are in the kitchen.
  \item \emph{Recall:} say the Punjabi for to grind, then the Punjabi for a wooden mortar, then say \emph{mohlā} again
  \item \emph{Write it:} copy \textbf{\pa{ਮੋਹਲਾ}} once, checking the whole word
\end{itemize}

\begin{tcolorbox}[breakable,colback=teal!4,colframe=teal!35!black,title={Before you move on}]
What does \pa{ਮੋਹਲਾ} mean? (\textquotedblleft{}A pestle\textquotedblright{}.) Say it once more.
\end{tcolorbox}
\section[\texorpdfstring{\pa{ਕੁੱਟਣਾ} (kuṭṭṇā)}{ਕੁੱਟਣਾ (kuṭṭṇā)}]{\pa{ਕੁੱਟਣਾ} (kuṭṭṇā) — to pound}

\label{lesson:PA-C240-kuttna}



\begin{quote}
Before the new one: say the Punjabi for a wooden mortar, then the Punjabi for a pestle.
\end{quote}

\subsection*{You'll want to know: \pa{ਕੁੱਟਣਾ}}
\textbf{\pa{ਕੁੱਟਣਾ}} — \emph{kuṭṭṇā} — \textquotedblleft{}to pound\textquotedblright{}.

\begin{quote}
\pa{ਮੈਂ} \pa{ਮੋਹਲੇ} \pa{ਨਾਲ} \pa{ਲਸਣ} \pa{ਕੁੱਟਦੀ} \pa{ਹਾਂ।}
\end{quote}

\begin{quote}
\emph{maiṁ mohle nāl lasaṇ kuṭṭdī hāṁ.}
\end{quote}

\begin{quote}
I pound garlic with the pestle.
\end{quote}

Past (I …): \pa{ਮੈਂ} \pa{ਅਦਰਕ} \pa{ਕੁੱਟਿਆ}, \emph{maiṁ adrak kuṭṭiā}.

\begin{cousinweb}
\textbf{\pa{ਕੁੱਟਣਾ}} (\emph{kuṭṭṇā}) is inherited from Sanskrit \emph{kuṭṭayati}, \textquotedblleft{}crushes, pounds\textquotedblright{}; CDIAL lists Punjabi \emph{kuṭṭṇā}, \textquotedblleft{}to beat, pound, ram\textquotedblright{} (CDIAL 3241).
\end{cousinweb}

\begin{grammarlens}[title={What you've built}]
One more word for this chapter.
\end{grammarlens}

\subsection*{Your turn}
\begin{itemize}
\raggedright
  \item \emph{Say it:} \emph{kuṭṭṇā}
  \item \emph{Say it:} \emph{kuṭṭṇā}, once more
  \item \emph{Say it:} \emph{maiṁ mohle nāl lasaṇ kuṭṭdī hāṁ.} — I pound garlic with the pestle.
  \item \emph{Recall:} say the Punjabi for a wooden mortar, then the Punjabi for a pestle, then say \emph{kuṭṭṇā} again
  \item \emph{Write it:} copy \textbf{\pa{ਕੁੱਟਣਾ}} once, checking the whole word
\end{itemize}

\begin{tcolorbox}[breakable,colback=teal!4,colframe=teal!35!black,title={Before you move on}]
What does \pa{ਕੁੱਟਣਾ} mean? (\textquotedblleft{}To pound\textquotedblright{}.) Say it once more.
\end{tcolorbox}
